\chapter{Substructuring and component reduction}\label{ch:cms}

\begin{tcolorbox}[colback=hpale,colframe=hnavy!60,coltitle=white,colbacktitle=hnavy!85,title={In this chapter}]
Real structures are rarely built, meshed and owned by one person: a wing, an engine and a pylon come from different teams, and each team can reduce its own part. The difficulty is joining the reduced parts again. If the reduction hides the interface, the parts cannot be connected; if it keeps every interface unknown, the interface can dominate the model. Component mode synthesis reduces each part separately while keeping the interface degrees of freedom physical, so the reduced parts couple like finite elements. We derive static (Guyan) condensation and the Craig--Bampton method, say what free-interface and interface reduction change, and describe the component-wise POD of non-matching nonlinear parts tied by a mortar operator. Each is measured: coupled halves of a beam against the full model, Guyan at equal size, the effect of reducing the interface of a membrane, what free-interface models do and do not carry, and component POD with a reused basis. The meaning of ``partitioning'' used here is the one deferred from Chapter~\ref{ch:mesh}.\par\smallskip
After it you should be able to build and compare Guyan and Craig--Bampton models, judge an assembly by its frequencies, and say what free-interface and interface reduction change.
\end{tcolorbox}

\section{The problem}\label{sec:cm-problem}

Take a clamped--clamped beam and cut it in the middle. Two teams each get one half and each reduces it. The first obvious reduction is to keep a few chosen nodes and let the rest follow statically, which is Guyan condensation; it is exact for static loads and it is easy. But how good is it for vibration, and can a better reduction be had for the same number of unknowns? The beam of 40 elements and its two halves are the running example of this chapter.

An assembly of $s$ substructures has stiffness and mass matrices $\bm K^{(k)},\bm M^{(k)}$ for each part $k$. Split the degrees of freedom of one part into those on the interface with other parts ($b$, \emph{boundary}) and the rest ($i$, \emph{internal}):
\begin{equation}\label{eq:cm-split}
\bm K=\begin{bmatrix}\bm K_{ii}&\bm K_{ib}\\\bm K_{bi}&\bm K_{bb}\end{bmatrix},\qquad
\bm M=\begin{bmatrix}\bm M_{ii}&\bm M_{ib}\\\bm M_{bi}&\bm M_{bb}\end{bmatrix}.
\end{equation}
The aim is a reduced model of each part that keeps $\bm u_b$ as coordinates and replaces the (many) internal unknowns by a few amplitudes, so that the parts can be joined by identifying interface unknowns. Whether a reduced model is good is judged on the assembly, by its natural frequencies against the full model's, not on the part alone.

\section{Mathematical structure}\label{sec:cm-math}

\begin{sidenote}{Interface, boundary, internal}
The unknowns where two parts touch are the interface or \emph{boundary} unknowns ($b$); everything else is \emph{internal} ($i$). Reductions of a part act on the internal unknowns and decide what to do with the interface. A method that keeps the boundary unknowns as physical displacements lets the parts be joined by simple addition.
\end{sidenote}

\begin{sidenote}{Static condensation}
Split the unknowns of a component into those that are kept (the interface or any chosen \emph{master} set) and those that are eliminated. Eliminating the second group by solving their equilibrium equations in terms of the first gives a smaller stiffness matrix, the \emph{Schur complement} $\bm K_{bb}-\bm K_{bi}\bm K_{ii}^{-1}\bm K_{ib}$, which is exact for static loads applied only to the kept unknowns. In dynamics the eliminated unknowns also have inertia, which condensation ignores, so the error grows with frequency.
\end{sidenote}

\paragraph{Guyan (static) condensation.}
Keep the interface (or any chosen master set) and let the others follow statically, $\bm K_{ii}\bm u_i+\bm K_{ib}\bm u_b=\bm 0$:
\begin{equation}\label{eq:cm-guyan}
\bm u=\bm T_G\bm u_b,\qquad \bm T_G=\begin{bmatrix}\bm I\\-\bm K_{ii}^{-1}\bm K_{ib}\end{bmatrix},\qquad \bm K_G=\bm T_G^T\bm K\bm T_G,\quad \bm M_G=\bm T_G^T\bm M\bm T_G .
\end{equation}
The condensed stiffness is exact for any static load on the master unknowns. The mass is approximate, because the internal motion is assumed static at all frequencies, so the method is accurate for modes much lower than the first mode of the internal unknowns with the interface held~\cite{guyan1965}.

\begin{sidenote}{Constraint modes and fixed-interface modes}
A constraint mode is the static shape of the part when one interface unknown is moved by one unit and the others are held. A fixed-interface mode is a vibration mode of the part with the interface clamped. The first kind carries the interface, the second the internal dynamics.
\end{sidenote}

\paragraph{Craig--Bampton.}
Add the missing dynamics of the internal unknowns as the vibration modes of the part with the interface fixed, $\bm K_{ii}\bm\Phi=\bm M_{ii}\bm\Phi\bm\Omega^2$ with $\bm\Phi^T\bm M_{ii}\bm\Phi=\bm I$, and keep the lowest $n_m$ of them~\cite{hurty1965,craig1968}:
\begin{equation}\label{eq:cm-cb}
\begin{bmatrix}\bm u_b\\\bm u_i\end{bmatrix}=\begin{bmatrix}\bm I&\bm 0\\\bm\Psi&\bm\Phi_{n_m}\end{bmatrix}\begin{bmatrix}\bm u_b\\\bm\eta\end{bmatrix},\qquad \bm\Psi=-\bm K_{ii}^{-1}\bm K_{ib}.
\end{equation}
The columns of $\bm\Psi$ are the constraint modes (the Guyan transformation), the columns of $\bm\Phi_{n_m}$ are the fixed-interface normal modes. With this basis the reduced stiffness is block diagonal, $\bm K_r=\mathrm{diag}(\bm K_G,\bm\Omega^2_{n_m})$, and the modal mass block is the identity. Because $\bm u_b$ are physical displacements, parts are joined by adding the reduced matrices at the shared interface unknowns, which enforces displacement compatibility. The reduced size of an assembly is the number of global interface unknowns plus $\sum_k n_m^{(k)}$. If all fixed-interface modes are kept the transformation spans the full space and the assembly is exact; with fewer, the error falls as $n_m$ grows because every discarded mode is above the retained ones.

\begin{sidenote}{Fixed-interface and free-interface modes}
The \emph{fixed-interface} modes of a component are its vibration modes with the interface held in place, and they describe how the inside moves relative to the interface. The \emph{free-interface} modes are those of the component floating free, which is closer to how it behaves when it is attached to a flexible neighbour. Truncating either set leaves out the static effect of the discarded modes; \emph{residual flexibility} is the correction that restores it.
\end{sidenote}

\paragraph{Free-interface reduction.}
Instead of fixing the interface, retain the lowest $n_m$ modes of the part with the interface free, $\bm\Phi_k$, and enrich them with the residual flexibility of the discarded modes at the interface, $\bm G_r=\bm K^{-1}-\bm\Phi_k\bm\Omega_k^{-2}\bm\Phi_k^T$, restricted to the interface columns. Rubin's method~\cite{rubin1975} keeps the inertia of these residual vectors; MacNeal's~\cite{macneal1971} neglects it and condenses the residual coordinates statically. The substructure must be restrained: free rigid-body motion is not handled.

\paragraph{Interface reduction.}
Craig--Bampton keeps every interface unknown, which is expensive when the interface is a surface with many nodes. Replace $\bm u_b$ by a few amplitudes of interface shapes, $\bm u_b=\bm\Psi_\Gamma\bm q_\Gamma$. Characteristic constraint modes~\cite{castanier2001} take the lowest eigenvectors of the interface block of the Craig--Bampton model, $\bm K_{bb}^{r}\bm x=\lambda\bm M_{bb}^{r}\bm x$ (with the reduced $bb$ blocks). Another choice is the Gram--Schmidt orthogonalisation of the interface rows of modes of the full model. Either gives a reduced model with $n_\Gamma+n_m$ coordinates; with all interface modes kept the model is unchanged.

\begin{algbox}{Interface reduction of a Craig--Bampton model}
\begin{pseudo}
\Require Craig--Bampton model with interface blocks $\bm K^r_{bb},\bm M^r_{bb}$, number of interface modes $n_\Gamma$
\State solve $\bm K^r_{bb}\bm x=\lambda\bm M^r_{bb}\bm x$;\ keep the $n_\Gamma$ lowest eigenvectors as $\bm\Psi_\Gamma$ \Comment{characteristic constraint modes; or Gram--Schmidt on mode rows}
\State $\bm T_\Gamma\gets\mathrm{diag}(\bm\Psi_\Gamma,\bm I_{n_m})$;\ project the reduced matrices
\Ensure model with $n_\Gamma+n_m$ coordinates per part set
\end{pseudo}
\end{algbox}

One eigenproblem of size $n_b$, the interface of the reduced model, is solved, which is cheap compared with the part. The reduction fails when the interface shapes that matter in the band are not among the lowest interface modes; with all of them kept the model is unchanged. The code is \anchor{rom_engine.interface_reduction.interface_reduced_craig_bampton} and \anchor{rom_engine.interface_reduction.gram_schmidt_interface_modes}; \Cref{fig:cm-iface} shows that the two interface bases coincide.

\tikzfigure[0.98\linewidth]{ch11_interface}{Interface reduction replaces the interface unknowns of a Craig--Bampton model by a few interface modes.}{fig:cm-ifaceidea}

\begin{sidenote}{Mortar coupling}
Where two components have different meshes along their common interface, the nodes do not match and cannot simply be identified. A \emph{mortar} method enforces continuity in a weak sense: the gap between the two displacement fields is required to be orthogonal to a set of test functions on the interface, with Lagrange multipliers as the unknown forces that tie them. It works for non-matching meshes at the price of the extra unknowns.
\end{sidenote}

\paragraph{Component POD with a mortar tie.}
For nonlinear parts, or parts with non-matching interface meshes, the part is reduced from nonlinear solutions computed with \emph{prescribed interface displacements}~\cite{ritzert2025}. The interface values are parametrised by a few smooth functions (translation, rotation, Legendre polynomials) whose coefficients, together with a load factor, are sampled with an optimal Latin hypercube, so that the basis of a part does not depend on the neighbour. In an assembly of a master part $a$ and a slave part $b$ tied by the mortar operator $\bm P$ (slave interface $=\bm P\times$ master interface), the unknowns are $\bm x=[\bm q_a;\bm q_b]$ with
\begin{equation}\label{eq:cm-comp}
\bm u_a=\bm V_a\bm q_a,\qquad \bm u_{b,\mathrm{int}}=\bm\Phi_b\bm q_b,\qquad \bm u_{b,\Gamma}=\bm P(\bm V_a\bm q_a)_{\Gamma_a},
\end{equation}
and the reduced Newton system is $\bm T^T\bm R(\bm T\bm x)=\bm 0$ with tangent $\bm T^T\bm K_t\bm T$. The forces are still evaluated on the full meshes (no hyper-reduction, Chapter~\ref{ch:hyper}), so the saving is in the linear algebra only.

\begin{algbox}{Component POD with a mortar tie}
\begin{pseudo}
\Require parts $a$ (master), $b$ (slave), mortar operator $\bm P$, interface parametrisation, sample count
\State sample interface-displacement coefficients and a load factor with an optimal Latin hypercube
\For{each sample, for each part}
  \State solve the nonlinear part problem with the \emph{prescribed} interface displacement;\ store the snapshot
\EndFor
\State POD of the snapshots of each part $\to\bm V_a,\ \bm\Phi_b$ \Comment{independent of the neighbour}
\State assemble $\bm T$ from \eqref{eq:cm-comp};\ \textbf{Newton:} $\bm T^T\bm R(\bm T\bm x)=\bm0$ with tangent $\bm T^T\bm K_t\bm T$
\Ensure reduced solution $\bm x$ and full displacement $\bm T\bm x$
\end{pseudo}
\end{algbox}

Offline there are samples $\times$ parts nonlinear solves; online the forces are still evaluated on the full meshes, so only the linear algebra is reduced. The method fails when the sampled interface functions do not span what the neighbour imposes, and the mortar tie is not exact for every choice of interface meshes. The code is \anchor{rom_engine.component_pod.ComponentBasis} and \anchor{rom_engine.component_pod.ReducedAssembly}; \Cref{tab:cm-comp} and \Cref{fig:cm-comp} are the checks.

\tikzfigure[0.98\linewidth]{ch11_cb}{Craig--Bampton substructuring: the interface unknowns are physical and shared, so parts are joined by adding reduced matrices.}{fig:cm-cbidea}

\begin{algbox}{Craig--Bampton substructuring of an assembly}
\begin{pseudo}
\Require parts $k$ with $\bm K^{(k)},\bm M^{(k)}$, interface lists, modes per part $n_m^{(k)}$
\State list the interface unknowns in the same physical order in every part that shares them
\For{each part $k$}
  \State solve $\bm K_{ii}\bm\Phi=\bm M_{ii}\bm\Phi\bm\Omega^2$ (mass-normalised);\ keep $n_m^{(k)}$ lowest modes
  \State $\bm\Psi\gets-\bm K_{ii}^{-1}\bm K_{ib}$;\ $\bm T\gets\begin{bmatrix}\bm I&\bm0\\\bm\Psi&\bm\Phi_{n_m}\end{bmatrix}$;\ $\bm K_r\gets\bm T^T\bm K\bm T$,\ $\bm M_r\gets\bm T^T\bm M\bm T$
\EndFor
\State add the reduced matrices at the global interface indices \Comment{displacement compatibility}
\State solve the reduced eigenproblem or response;\ recover each part's displacement with its own $\bm T$
\Ensure reduced assembly of size (global interface unknowns) $+\sum_kn_m^{(k)}$
\end{pseudo}
\end{algbox}

For each part there is one eigenproblem of the interior block and one factorisation of $\bm K_{ii}$ with $n_b$ back-substitutions. The interface unknowns are all kept, so a large interface is expensive; a part with rigid-body motion when the interface is fixed does not occur, and too few retained modes show as an error that falls as $n_m$ grows. The code is \anchor{rom_engine.cms.craig_bampton}, \anchor{rom_engine.cms.couple} and \anchor{rom_engine.cms.guyan}; \Cref{fig:cm-cb} gives the error against the reduced size, with Guyan as the comparison.

\section{Implementation in the packages}\label{sec:cm-impl}

\begin{lstlisting}[caption={Two Craig--Bampton halves coupled at two interface dofs (KL, ML, ibL: left half and its interface dof positions; KR, MR, ibR: right half).},label={lst:cm-run}]
from rom_engine.cms import craig_bampton, couple, guyan

cl = craig_bampton(KL, ML, ibL, n_modes=5)        # constraint modes + 5 fixed-interface modes
cr = craig_bampton(KR, MR, ibR, n_modes=5)
cp = couple([cl, cr], [[0, 1], [0, 1]])            # global interface indices of each part
freq_hz, shapes = cp.solve_modal(4)
u_left = cp.expand(0, shapes[:, 0])               # full-order displacement of part 0, mode 1

g = guyan(K, M, master_dofs)                       # static condensation of one model
\end{lstlisting}

\paragraph{Try it.} \Cref{lst:cm-try} puts the two reductions side by side on the beam, at the same size. Guyan condenses the whole beam on twelve evenly spaced transverse masters; Craig--Bampton couples the two halves with five fixed-interface modes each, which also gives twelve unknowns (the two interface unknowns plus the retained modes of both parts).

\begin{lstlisting}[caption={Guyan condensation against two coupled Craig--Bampton halves, both of 12 unknowns (run from the repository root with both packages on the path).},label={lst:cm-try}]
import numpy as np
from scipy.linalg import eigh
from fea_engine import elements
from fea_engine.geometry import generate_mesh
from fea_engine.material import EI_beam, Material, Section
from fea_engine.solver import FESystem
from rom_engine.cms import couple, craig_bampton, guyan

MAT, SEC, N = Material(E=210e9, nu=0.3, rho=7800.0), Section(A=0.01, I=8.33e-6), 20

def beam(L, n, fixed):                       # free-dof matrices of a clamped beam piece
    s = FESystem(generate_mesh(dim=1, L=L, n=n), elements.Beam2DEulerBernoulli())
    s.assemble_stiffness(EI_beam(MAT, SEC)); s.assemble_mass(7800.0 * 0.01); s.fix_dofs(fixed, [0, 1])
    fr = np.asarray(s.free_dofs)
    return fr, np.array(s.K)[np.ix_(fr, fr)], np.array(s.M)[np.ix_(fr, fr)]

hz = lambda K, M: np.sqrt(eigh(K, M, eigvals_only=True)[:4]) / (2 * np.pi)
fr, K, M = beam(1.0, 2 * N, [0, 2 * N])      # the whole beam: the reference
f_full = hz(K, M)

# naive: Guyan condensation on 12 evenly spaced transverse masters
v = np.arange(0, K.shape[0], 2)
masters = v[np.round(np.linspace(0, len(v) - 1, 14)[1:-1]).astype(int)]
g = guyan(K, M, masters)
err_g = np.max(np.abs(hz(g.K, g.M) - f_full) / f_full)

# Craig-Bampton: two halves, 5 fixed-interface modes each, joined at the shared node
frL, KL, ML = beam(0.5, N, [0]); frR, KR, MR = beam(0.5, N, [N])
pos = lambda f, d: int(np.where(f == d)[0][0])
cl = craig_bampton(KL, ML, [pos(frL, 2 * N), pos(frL, 2 * N + 1)], 5)
cr = craig_bampton(KR, MR, [pos(frR, 0), pos(frR, 1)], 5)
cp = couple([cl, cr], [[0, 1], [0, 1]])
err_c = np.max(np.abs(cp.solve_modal(4)[0] - f_full) / f_full)

print(f"full beam: {K.shape[0]} unknowns, f1 = {f_full[0]:.2f} Hz")
print(f"Guyan           {g.K.shape[0]:2d} unknowns: worst error of the first 4 frequencies {err_g:.1e}")
print(f"Craig-Bampton   {cp.K.shape[0]:2d} unknowns: worst error of the first 4 frequencies {err_c:.1e}")
\end{lstlisting}

\begin{lstlisting}[numbers=none,language={},keywordstyle={},commentstyle={},frame=single,aboveskip=2pt]
full beam: 78 unknowns, f1 = 533.25 Hz
Guyan           12 unknowns: worst error of the first 4 frequencies 1.1e-03
Craig-Bampton   12 unknowns: worst error of the first 4 frequencies 9.3e-05
\end{lstlisting}

At the same size Craig--Bampton is more than ten times more accurate. Guyan assumes the interior moves statically at every frequency, which is only right for the lowest modes; the fixed-interface modes put the interior dynamics back. The next section follows the error as the number of modes grows.

\begin{table}[ht]
\centering\small
\begin{tabularx}{\linewidth}{@{}l X >{\centering\arraybackslash}p{0.11\linewidth}@{}}
\toprule
Concept & Code & Script calls \\
\midrule
static (Guyan) condensation & \anchor{rom_engine.cms.guyan} & yes \\
Craig--Bampton substructure and coupling & \anchor{rom_engine.cms.craig_bampton}, \anchor{rom_engine.cms.couple} & yes \\
free-interface reduction (Rubin, MacNeal) & \anchor{rom_engine.cms.free_interface_cms} & yes \\
interface reduction by characteristic constraint modes & \anchor{rom_engine.cms.reduce_interface} & yes \\
interface modes from other mode sets & \anchor{rom_engine.interface_reduction.gram_schmidt_interface_modes} & yes \\
interface-reduced Craig--Bampton & \anchor{rom_engine.interface_reduction.interface_reduced_craig_bampton} & yes \\
frequency error metric & \anchor{rom_engine.interface_reduction.natural_frequency_errors} & yes \\
MAC metric & \anchor{rom_engine.interface_reduction.mac_diagonal} & yes \\
component basis from snapshots & \anchor{rom_engine.component_pod.ComponentBasis} & yes \\
mortar-tied reduced assembly & \anchor{rom_engine.component_pod.ReducedAssembly} & yes \\
mortar tie operator (in \fea{}) & \anchor{fea_engine.mortar.MortarTie2D} & yes \\
worked experiment & \path{paper/experiments/e3_cms.py} & -- \\
\bottomrule
\end{tabularx}
\caption{Where the concepts of this chapter live in the repository (tag \booktag). The dual Craig--Bampton, joint-interface-mode and Fourier interface bases of \code{interface_reduction.py} are not tested here. The last column says whether this chapter's own checks call, by name, the code listed in the row: yes (every name), part, or no. Code reached only through another routine counts as no, so the column is a lower bound on what this book exercises.}\label{tab:cm-map}
\end{table}

\section{Verification}\label{sec:cm-verif}

All numbers were measured on \machine{}. The package never imports \fea{} (Chapter~\ref{ch:arch}), so every experiment passes plain arrays.

\subsection{Coupled halves of a beam, and Guyan at equal size}

A clamped--clamped beam of $40$ Euler--Bernoulli elements (\ccmBeamFull{} free unknowns; $f_1=\ccmBeamFOne$~Hz) is split in two halves of $20$ elements, each clamped at its outer end, with the two transverse and rotation unknowns of the shared node as the interface (model of \path{paper/experiments/e3_cms.py}). The reference is the generalised eigenproblem of the unreduced beam. \Cref{fig:cm-cb} shows the largest relative error of the first four frequencies. With one fixed-interface mode per half (a \ccmCbOneSize{}-unknown model) the error is \ccmCbOneErr; with five (\ccmCbFiveSize{} unknowns) \ccmCbFiveErr; with twenty (\ccmCbTwentySize) \ccmCbTwentyErr; and with all of them (\ccmCbAllSize{} unknowns, the full space) \ccmCbAllErr, which is the exactness limit to rounding. The error decreases monotonically at every size tested (\ccmCbMonotone).

At equal size Guyan condensation is much less accurate: with $12$ transverse master unknowns (the size of the $n_m=5$ Craig--Bampton model) the error is \ccmGuyTwelveErr, about a factor $12$ larger than \ccmCbFiveErr, and Guyan needs about $20$ masters to reach \ccmGuyTwentyErr, a level that Craig--Bampton already beats at $12$ unknowns (\ccmCbFiveErr).  Guyan is accurate where it should be: its first frequency, with $12$ masters, has an error of \ccmGuyOneErrFirst, and the error is dominated by the fourth mode (\ccmGuyTwelveHigh). It is also exact statically: the displacement of the full beam under a midspan load is reproduced by a single master at the load point to \ccmGuyStatic{} in relative norm.

\begin{figure}[ht]
\centering
\includegraphics[width=0.62\linewidth]{figures/ch11_cb.pdf}
\caption{Largest relative error of the first four frequencies against the reduced size, Craig--Bampton (two halves) and Guyan (transverse masters, evenly spaced).}\label{fig:cm-cb}
\end{figure}

\subsection{Reducing the interface}

A membrane (a scalar field on a $16\times21$ grid, \ccmMemN{} unknowns, left edge grounded, right edge as the interface of \ccmMemIface{} unknowns) is reduced standalone and compared with its own full model over the first $15$ frequencies. \Cref{fig:cm-iface} (left) shows the Craig--Bampton error against the size: \ccmMemNmTenErr{} with $n_m=10$ ($\ccmMemNmTenSize$ unknowns) and \ccmMemNmFortyErr{} with $n_m=40$ ($\ccmMemNmFortySize$ unknowns). The right panel fixes $n_m=40$ and reduces the \ccmMemIface{} interface coordinates to $k$ amplitudes of characteristic constraint modes, or of interface modes built from modes of the full model. With $2$ interface coordinates the error is \ccmCcTwoErr, with $4$ it is \ccmCcFourErr, and from $6$ coordinates upwards the error equals that of the unreduced interface (\ccmPlainErr), so $15$ of the $21$ interface coordinates were redundant for these $15$ modes. The mode-shape agreement (the smallest diagonal MAC over the $15$ modes) is \ccmCcFourMac{} at $4$ coordinates and \ccmCcEightMac{} at $8$. The two interface bases gave the same errors at every $k$ in this test, and the subspaces they span were found to coincide to $10^{-14}$; we did not investigate whether this holds beyond this membrane.

\begin{figure}[ht]
\centering
\includegraphics[width=0.98\linewidth]{figures/ch11_interface.pdf}
\caption{Left: Craig--Bampton error of the standalone membrane against reduced size. Right: error with $40$ fixed-interface modes against the number of interface coordinates kept; the two interface bases coincide.}\label{fig:cm-iface}
\end{figure}

\subsection{Free-interface reduction}

\Cref{tab:cm-free} compares modal truncation, Craig--Bampton, Rubin and MacNeal for a spring--mass chain of $40$ masses with the interface at the free end and the middle node. These models are standalone: the package couples only Craig--Bampton models (\code{couple} needs physical interface unknowns), so this tests what each model carries and not how it behaves in an assembly.

\begin{table}[ht]
\centering\small
\begin{tabular}{rlrrrr}
\toprule
$n_m$ & model & size & lowest $n_m$ freq. & modes $n_m{+}1$ to size & static tip error \\
\midrule
3 & modal truncation & 3 & 0 & -- & $8.5\times10^{-1}$ \\
 & Craig--Bampton & 5 & $1.4\times10^{-2}$ & $8.5\times10^{-1}$ & $1.4\times10^{-16}$ \\
 & Rubin & 5 & $3.2\times10^{-14}$ & $8.4\times10^{-1}$ & $1.4\times10^{-16}$ \\
 & MacNeal & 3 & $3.2\times10^{-14}$ & -- & $8.5\times10^{-1}$ \\
\midrule
6 & modal truncation & 6 & 0 & -- & $7.1\times10^{-1}$ \\
 & Craig--Bampton & 8 & $1.6\times10^{-2}$ & $6.3\times10^{-1}$ & $1.4\times10^{-16}$ \\
 & Rubin & 8 & $3.2\times10^{-14}$ & $6.7\times10^{-1}$ & $1.4\times10^{-16}$ \\
 & MacNeal & 6 & $3.2\times10^{-14}$ & -- & $7.1\times10^{-1}$ \\
\midrule
10 & modal truncation & 10 & 0 & -- & $5.3\times10^{-1}$ \\
 & Craig--Bampton & 12 & $1.6\times10^{-2}$ & $3.4\times10^{-1}$ & $1.4\times10^{-16}$ \\
 & Rubin & 12 & $3.2\times10^{-14}$ & $4.4\times10^{-1}$ & $1.4\times10^{-16}$ \\
 & MacNeal & 10 & $3.2\times10^{-14}$ & -- & $5.3\times10^{-1}$ \\
\bottomrule
\end{tabular}

\caption{Free-interface and Craig--Bampton models of a chain of $40$ masses, interface at two nodes. ``lowest $n_m$ freq.'': largest relative error of the lowest $n_m$ frequencies; ``modes $n_m{+}1$ to size'': the frequencies above those, up to the model size; ``static tip error'': relative error of the displacement at the end node under a unit force at that node.}\label{tab:cm-free}
\end{table}

Two results. The retained frequencies of the free-interface models are reproduced to \ccmFreeRubinLow{} because the retained modes are in the span of the basis by construction; this is not an accuracy result, and Craig--Bampton has an error of \ccmFreeCbLowThree{} on the same modes at $n_m=3$. The static flexibility is the real difference: Craig--Bampton and Rubin reproduce the static response at the interface node to rounding (\ccmFreeCbStatic{} and \ccmFreeRubinStatic{}), while modal truncation misses it by \ccmFreeModalStatic{} at $n_m=3$ (\ccmFreeModalStaticTen{} at $n_m=10$). The \texttt{macneal} model, as returned by the package, has the same static error as modal truncation (\ccmFreeMacStatic{} and \ccmFreeMacStaticTen{}) and the same frequencies: the residual coordinates are decoupled from the modal ones in the stiffness, so condensing them out leaves $\bm K_r=\bm\Omega_k^2$ and the modal mass, and the residual flexibility is not carried by the returned $(\bm T,\bm K,\bm M)$ for a static load. We report this as observed behaviour, not as a defect; the package's test file for these models checks the frequencies only. The frequencies above the retained ones are poorly predicted by every model with a handful of states (errors of $0.3$ to $0.9$).

\subsection{Component POD with a mortar tie}

Two nonlinear scalar membranes (a cubic hardening term $4u^3$, clamped outer edges), a master of $9\times9$ nodes and a slave of $9\times7$ nodes with a non-matching interface, are tied by a mortar operator from \code{fea_engine.mortar}; the full tied problem has \ccmCompFull{} unknowns and is solved by Newton. Each component basis comes from $30$ snapshots with Legendre-parametrised interface values and a load factor, plus the zero and axis points. \Cref{tab:cm-comp} gives the error of the reduced assembly against the full tied solution: \ccmCompErrThree{} with $3$ modes per part ($6$ unknowns), \ccmCompErrSix{} with $6$ and \ccmCompErrTwelve{} with $12$ (\ccmCompRedTwelve{} unknowns of \ccmCompFull{}); the error decreases with the basis size (\ccmCompMono), Newton converged in all cases (\ccmCompConv), and the slave interface obeys the mortar tie to \ccmCompTie. The master basis was then reused against a different slave component, a coarser $9\times5$ mesh with its own slave basis: the reduced solution agrees with the full tied solution of that new assembly to \ccmCompReuseErr{} (converged: \ccmCompReuseConv), without recomputing the master snapshots.

\begin{table}[ht]
\centering\small
\begin{tabular}{rrrrrr}
\toprule
rank per component & reduced unknowns & full unknowns & relative error & Newton iterations & tie violation \\
\midrule
3 & 6 & 144 & $5.3\times10^{-2}$ & 8 & $1.1\times10^{-16}$ \\
6 & 12 & 144 & $2.1\times10^{-3}$ & 8 & $2.2\times10^{-16}$ \\
12 & 24 & 144 & $2.8\times10^{-4}$ & 9 & $2.2\times10^{-16}$ \\
\bottomrule
\end{tabular}

\caption{Component POD of two non-matching nonlinear membranes tied by a mortar operator: error of the reduced assembly against the full tied solution.}\label{tab:cm-comp}
\end{table}

\begin{figure}[ht]
\centering
\includegraphics[width=0.55\linewidth]{figures/ch11_component.pdf}
\caption{Relative error of the component-POD assembly against the POD rank of each component.}\label{fig:cm-comp}
\end{figure}

\section{Limits and failure modes}\label{sec:cm-limits}

\begin{limitsbox}
\begin{itemize}
\item \textbf{Small, simple test models.} A beam, a graph-Laplacian membrane, a spring chain and two scalar nonlinear membranes: no three-dimensional solid substructure, no shell, no assembly of more than two parts, no rotating or cyclic structure. The Fourier and multistage bases, the dual Craig--Bampton and the joint-interface-mode method are in the package but not tested here.
\item \textbf{Fixed-interface modes need a positive definite $\bm K_{ii}$}; the package raises an error if the interface does not constrain the internal unknowns.
\item \textbf{Free-interface models are not coupled by the package.} The \texttt{macneal} model does not carry the residual flexibility in its returned matrices (\Cref{tab:cm-free}), and the frequency test of a standalone model is satisfied by construction.
\item \textbf{Interface reduction was tested on one membrane}, for the first $15$ modes, where six coordinates sufficed. The number needed depends on the frequency range and the interface, no rule for choosing it is given, and the agreement of the two interface bases was observed and not explained.
\item \textbf{Component POD} evaluates forces on the full meshes, so the saving is in the linear algebra only; it covers small displacements through a constant mortar operator, two components per call, and a scalar membrane. Hyper-reduction of the component forces is the subject of Chapter~\ref{ch:hyper}, and timing was not measured.
\end{itemize}
Guyan is accurate for low modes and statics only (\Cref{fig:cm-cb}).
\end{limitsbox}

\section{Summary}\label{sec:cm-summary}

We started by cutting a beam in two and asking how to reduce the halves so that they could be joined again. The answer is to keep the interface physical. Guyan condensation does this statically and is exact for static loads, but it loses accuracy in the higher modes at the same size. Craig--Bampton adds the fixed-interface modes of each part to the same constraint modes, and in the test it was accurate, improved steadily as modes were added, and became exact when all of them were kept.

Two further questions came with larger interfaces and nonlinear parts. A few interface modes were enough for the first modes of the test membrane, and free-interface models contain their retained modes by construction, but whether they also carry the static flexibility depends on the variant: the package's MacNeal model does not, in the form it returns. Component POD gave an accurate mortar-tied reduced assembly of nonlinear parts, and one basis could be reused against a different neighbour. The next chapter drops the requirement that the matrices are available at all.

\section*{Exercises}\label{sec:cm-exercises}
\begin{enumerate}
\item \emph{(Implementation)} Split the beam of \Cref{lst:cm-run} into three parts with two interfaces. Does the monotone convergence in $n_m$ still hold, and what is the size of the global interface?
\item \emph{(Implementation)} Repeat \Cref{fig:cm-iface} (right) for the first $40$ frequencies instead of $15$. How many interface coordinates are redundant now?
\item \emph{(Interpretation)} Compute the static response of the chain of \Cref{tab:cm-free} to a force at the middle node, not at the end. Which models are still exact?
\item \emph{(Implementation)} Replace the cubic term of the component membranes by $8u^3$. How does the error of \Cref{tab:cm-comp} change at rank $12$, and does Newton still converge?
\item \emph{(Implementation)} Time the reduced Newton solve of \Cref{tab:cm-comp} against the full tied solve. Where does the time go, given that the forces are evaluated on the full meshes? (Test file: \texttt{rom\_engine/tests/test\_component\_pod.py}.)
\item \emph{(Derivation)} Show that if all fixed-interface modes are kept, the Craig--Bampton transformation $\left[\begin{smallmatrix}\bm I&\bm 0\\\bm\Psi&\bm\Phi\end{smallmatrix}\right]$ is invertible, and hence that the coupled model is exact.
\end{enumerate}

\chapterrefs
